\documentclass[10pt,a4paper]{amsart}
\usepackage[margin=19mm]{geometry}
\usepackage{amsmath,amssymb,mathtools}
\usepackage[scr=rsfs,cal=euler]{mathalfa}
\usepackage{xfrac}
\usepackage[unicode,hidelinks,bookmarks=false]{hyperref}
\newcommand{\ie}{i.e.\ }
\newcommand{\cf}{cf.\ }
\newcommand{\resp}{respectively\ }
\newcommand{\Subsection}{\subsection}
\providecommand{\nobreakdash}{\nobreak\hbox{-}}
\setlength{\emergencystretch}{2em}
\multlinegap0pt
\usepackage{bm}
\usepackage{bbm}
\usepackage{dsfont}
\usepackage{xspace}
\usepackage{cancel}
\usepackage{mathtools}
\usepackage{amsthm}
\makeatletter
\@ifundefined{lem}{\newtheorem{lem}{\bf Lemma}}{}
\makeatother
\title[The Cartesian product of exact approximation sets]{The Cartesian product of exact approximation sets}
\author{Jing Guo, Mumtaz Hussain, Bixuan Li,, Johannes Schleischitz}
\date{Source: arXiv:2608.26818v1 (CC BY 4.0)}
\begin{document}
\maketitle

\noindent\emph{Source and licence.} The Cartesian product of exact approximation sets. Version arXiv:2608.26818v1 (CC BY 4.0), distributed under the Creative Commons Attribution 4.0 International licence, as recorded in the arXiv licence metadata for this exact version and shown on the abstract page https://arxiv.org/abs/2608.26818v1. Full rights notice: licenses/arxiv-2608.26818-CC-BY-4.0.txt.\par\smallskip

\noindent\textit{Editorial scope.} Counted unit: the Lem whose source label is \texttt{count} in the article (source0.tex lines 899--901, source label \texttt{count}), delivered together with the article's own complete original proof of it (source0.tex lines 902--980, one \texttt{proof} environment of 79 lines). No mathematics is written, repaired or re-derived: statement and proof are the article's own text line for line. Required context retained uncounted: condithm1 (lines 349--352); WDP (lines 795--810); lcondi (lines 893--895). Every definition, display and label of those lines is retained unchanged. Omitted from this package and not redistributed: 1--348, 353--794, 811--892, 896--898, 981--1599. Nothing inside a retained excerpt is omitted or altered: every retained body is delivered exactly as the article's lines are written, so no omission receipt of any kind accompanies this package. Citations: the retained excerpts carry no citation command at all, so no bibliography entry of the article is transcribed and no reference is redistributed. Reference adaptation: none was needed and none was made; every reference inside the delivered unit resolves inside it, so no unresolved reference or citation is produced. Printed slips kept literal: no printed word, symbol, hypothesis or number of the counted statement or of its proof was altered. Layout and class: the article's own class is \texttt{amsart} with packages including geometry, amsmath, amsthm, xcolor, amssymb, graphicx, soul, ulem, relsize, enumerate, color, euscript, tikz; the wrapper is the lane's pinned \texttt{amsart} editorial wrapper at 10pt on A4, which restates the theorem-like environments the retained text needs and adds the article's own macro definitions that the retained text needs (none), with extra wrapper packages (bm, bbm, dsfont, xspace, cancel, mathtools). The delivered numbering is the unit's own. Assets: no retained span contains a figure environment, a graphics-inclusion command, a PDF-page inclusion, a table or a TikZ picture; no image file is copied into this package, the package declares no asset dependency and no image file is redistributed with it. Licence: version arXiv:2608.26818v1 of this article is distributed under the Creative Commons Attribution 4.0 International licence, as recorded in the arXiv licence metadata for this exact version and shown on the abstract page https://arxiv.org/abs/2608.26818v1. A full rights notice is delivered at licenses/arxiv-2608.26818-CC-BY-4.0.txt. Characters: every retained excerpt is pure ASCII, so the delivered engine, XeLaTeX with its default Latin Modern text font, reports no missing character.
\begin{align}\label{condithm1}
\psi_i(q)=o(q^{-2}),
\qquad i=1,\dots,m,
\end{align}

\begin{lem}\label{WDP}For any interval $U$, let $Q$ be sufficiently large such that
        \begin{equation}\label{f1}
        Q^2\cdot|U|\ge 16,\qquad
       {8Q^{-1}\log (QN^{-1})}\le \frac{1}{16}|U|.
        \end{equation} Then, there exists a subset $\mathcal{A}_{Q}(U)\subset \mathcal{Q}_{Q}$ such that
\begin{itemize}
  \item For any two distinct rationals $p/q$ and $p^{\prime}/q^{\prime}$ in $\mathcal{Q}_Q$, one has
   $$\left|\frac{p}{q}-\frac{p^{\prime}}{q^{\prime}}\right|\ge \frac{1}{Q^2}.$$
     \item The union of the following disjoint balls is  contained in $U$
        $$\bigcup_{\frac{p}{q}\in \mathcal{A}_{Q}(U)}B\left(\frac{p}{q},\frac{1}{2Q^2}\right)\subset U.$$
    \item The cardinality of \(\mathcal A_Q(U)\) satisfies
        \begin{align}\label{number}
        \# \mathcal{A}_{Q}(U)\asymp_{N} \mathcal{L}(U)\cdot Q^{2}.
        \end{align}
\end{itemize}
\end{lem}

 \begin{align}\label{lcondi}
    \sum_{1\le t\le k_-}N^{-2t}Q_+^2\psi(N^{-t}Q_+^2)\le \frac{1}{1152}(1-c)N^{-7}.
    \end{align}

\begin{lem}[Counting Lemma]\label{count}
    Under the notations as above, we have $\#\mathcal{D}_{Q_+}\le \#\mathcal{A}_{Q_+} / 2$.
\end{lem}

\begin{proof}\label{trequire}
    It readily follows, by the assumption {(\ref{condithm1})} on $\psi$, that
    \begin{align}\label{psicondi}
    \psi(q)<\frac{1}{2N^2}\cdot\frac{1}{q^2}\text{ for large }q.
    \end{align}
     We partition $\mathcal{D}_{Q_+}$ into finitely many subsets as follows. For any $1\le t\le k_-$, let
     \begin{align*}
\mathcal{D}_{Q_+,t}:= \Bigg\{ \frac{p'}{q'} \in \mathcal{A}_{Q_+} : & A_{c}(p'/q') \cap B\left(\frac{p}{q}, c\psi(q)\right) \neq \emptyset, \\
& \quad \text{for some rational point } \frac{p}{q} \text{ with } N^{-t}Q_+\le q< N^{-t+1}Q_+ \Bigg\}.
\end{align*}
Now, we use a volume argument to estimate its cardinality. For any rational $p^{\prime}/q^{\prime}\in\mathcal{D}_{Q_+,t}$, there exists a rational point $p/q$ such that $$ N^{-t}Q_+\le q<N^{-t+1}Q_+ \textrm{ and } \left|\frac{p}{q}-\frac{p^{\prime}}{q^{\prime}}\right|\le c\psi(q)+\psi(q^{\prime}).$$
We claim that ${p}/{q}\neq{p^{\prime}}/{q^{\prime}}.$ Otherwise, since $\textrm{gcd}(p^{\prime},q^{\prime})=1$, one has $q=rq^{\prime}$ for some integer $r\ge 1$. It follows that $q\ge q^{\prime}$. By the non-increasing of $\psi$, one has $c\psi(q)\le c\psi(q^{\prime})$. Therefore, by the definition of annulus, it follows that
$$A_{c}(p'/q') \cap B\left(\frac{p^{\prime}}{q^{\prime}}, c\psi(q^{\prime})\right) =\emptyset\ \Longrightarrow\  A_{c}(p'/q') \cap B\left(\frac{p}{q}, c\psi(q)\right) =\emptyset,$$
which yields a contradiction. On the other hand, for any two distinct rational points, one has
\begin{align*}
\left|\frac{p}{q}-\frac{p^{\prime}}{q^{\prime}}\right|\ge\frac{1}{qq^{\prime}}.
\end{align*}
By (\ref{psicondi}) applied to $q^{\prime}$, it follows that
$$\frac{1}{qq^{\prime}}\le  \left|\frac{p}{q}-\frac{p^{\prime}}{q^{\prime}}\right|\le  c\psi(q)+\psi(q^{\prime})\le c\psi(q)+\frac{1}{2N^2q^{\prime2}}.$$
Together with the range of $q$ and $q^{\prime}$ and $t\ge 1$, one has $$qq^{\prime}< N^{-t+1}Q_+\cdot Q_+\le Q_{+}^2
= N^2\cdot (N^{-1}Q_+)^2\le N^2q^{\prime2}, $$
where the second inequality holds since $t\ge 1$.
%{\color{blue} can you explain the last two implications a bit more? how do you use $\mathcal{A}_{Q_+}$, i dont get the argument }
    Combining the last two properties, we conclude that
    $$\frac{1}{2N^2}\cdot\frac{1}{q^{\prime2}}\le
    \frac{1}{qq^{\prime}}-\frac{1}{2N^2q^{\prime2}}
    \le c\psi(q),$$
    and consequently
    \begin{align*}
    \left|\frac{p}{q}-\frac{p^{\prime}}{q^{\prime}}\right|
    \le c\psi(q)+\psi(q^{\prime})
    \le 
    c\psi(q)+\frac{1}{2N^2q^{\prime2}}
    \le 2c\psi(q).
    \end{align*}
    Hence, using again $1/(2N^2q^{\prime2})\le c\psi(q)$, for any $p^{\prime}/q^{\prime}\in\mathcal{D}_{Q_+,t}$, we have found a rational point $p/q$ with $N^{-t} Q_+\le q<N^{-t+1}Q_+$ such that
    $$B\left(\frac{p^{\prime}}{q^{\prime}},\frac{1}{2N^2q^{\prime2}}\right)\subset
    B\left(\frac{p^{\prime}}{q^{\prime}},c\psi(q)\right)
    \subset B\left(\frac{p}{q},3c\psi(q)\right).$$
    Thus, we obtain the inclusion
    \begin{align}\label{result1}
       \bigcup_{p^{\prime}/q^{\prime}\in\mathcal{D}_{Q_+}}B\left(\frac{p^{\prime}}{q^{\prime}},\frac{1}{2N^2q^{\prime2}}\right)\subset \bigcup_{1\le t\le k_-}\bigcup_{\substack{p/q:N^{-t} Q_+\le q<N^{-t+1}Q_+}}B\left(\frac{p}{q},3c\psi(q)\right),
    \end{align}
    where the balls in the left-hand side are pairwise disjoint by the separation of rationals in $\mathcal{D}_{Q_+,t}$. We note that the second union of balls in the right-hand side should be further restricted to
    $$B\left(\frac{p}{q},c\psi(q)\right)\cap U_0\neq\emptyset.$$
    Thus, it leads us to estimate the cardinality of following set
    $$\mathcal{N}_{Q_+,t}=\left\{\frac{p}{q}:B\left(\frac{p}{q},c\psi(q)\right)\cap U_0\neq\emptyset, \ N^{-t}Q_+\le q< N^{-t+1}Q_+\right\}.$$
    By similar argument as above, for any $p/q\in\mathcal{N}_{Q_+,t}$, we have $p/q\neq p_0/q_0.$ Similarly, if $p/q= p_0/q_0$, then $q\ge q_0$ by $\textrm{gcd}(p_0,q_0)=1$. Since $U_0$ is one of the interval of $A_{c}(p_0/q_0)$, it follows that
    $$B\left(\frac{p_0}{q_0},c\psi(q_0)\right)\cap U_0=\emptyset\ \Longrightarrow\ B\left(\frac{p}{q},c\psi(q)\right)\cap U_0=\emptyset.$$
    %This is why we choose $A_{c}(p_0/q_0)$ as the starting rather than $A_{c_0}(p_0/q_0)$.
    Therefore, we have that
    $$\frac{1}{q^{2}}\le\frac{1}{q_0q}\le \left|\frac{p_0}{q_0}-\frac{p}{q}\right|\le c\psi(q)+\psi(q_0)<\frac{1}{2q^{2}}+\psi(q_0)\Longrightarrow \frac{1}{2q^{2}}\le \psi(q_0),$$
    where the fourth inequality holds by (\ref{psicondi}) with $N\ge 8$ and $c<1$.
    It implies that
    $$\left|\frac{p_0}{q_0}-\frac{p}{q}\right|\le 
    \frac{1}{2q^{2}}+\psi(q_0)\le 2\psi(q_0).$$
    %$$B\left(\frac{p}{q},\frac{c}{2N^2q^{2}}\right)\subset B\left(\frac{p_0}{q_0},3\psi(q_0)\right).$$
    %{\color{red} Factor $2$ not needed here?!}
    Then, for any rational point $p/q\in\mathcal{N}_{Q_+,t}$, we have that
    $$B\left(\frac{p}{q},\frac{1}{2N^2q^2}\right)\subset B\left(\frac{p_0}{q_0},3\psi(q_0)\right).$$
    Therefore, we have the following inclusion
    %by triangle inequality, for any $x\in B({p}/{q},{c}/{(2N^2q^{2}}))$, we have
    %$$\left|x-\frac{p_0}{q_0}\right|\le \left|x-\frac{p}{q}\right|+\left|\frac{p}{q}-\frac{p_0}{q_0}\right|\le \frac{1}{2q^{2}}+2\psi(q_0)\le 3\psi(q_0).$$
    $$\bigcup_{{p}/{q}\in\mathcal{N}_{Q_+,t}}B\left(\frac{p}{q},\frac{1}{2N^2q^{2}}\right)\subset B\left(\frac{p_0}{q_0},3\psi(q_0)\right).$$
    where the balls in the left-hand side are pairwise disjoint by separation of rationals in $\mathcal{N}_{Q_+,t}$. By a volume argument, one has
    $$\#\mathcal{N}_{Q_+,t}\cdot \frac{1}{2N^2\cdot (N^{-t+1}Q_+)^2}\le 3\psi(q_0).$$
    Returning back to (\ref{result1}), by a volume argument again, one has
    \begin{align*}
         \#\mathcal{D}_{Q_+}\cdot\frac{1}{2N^2\cdot Q_+^2}&\le \sum_{1\le t\le k_-}\#\mathcal{N}_{Q_+,t}\cdot 3c\psi(N^{-t}Q_+)\\
         &\le \sum_{1\le t\le k_-}{2N^{-2t+4}Q_+^2}\cdot3\psi(q_0)\cdot3\psi(N^{-t}Q_+).
    \end{align*}
    Thus, we conclude that
    \begin{align*}
    \#\mathcal{D}_{Q_+}&\le{36 \cdot Q_+^2 N^6}\psi(q_0)\cdot\sum_{1\le t\le k_-}N^{-2t}Q_+^2\psi(N^{-t}Q_+)\\
    &\le\#\mathcal{A}_{Q_+}\cdot 576(1-c)^{-1}N^7\sum_{1\le t\le k_-}N^{-2t}Q_+^2\psi(N^{-t}Q_+)\\
    &\le \#\mathcal{A}_{Q_+}/2,
    \end{align*}
    where the second inequality holds by the lower bound estimate of the third item in Lemma \ref{WDP}, the last inequality holds by the choice of $Q_-$ and $Q_+$ in (\ref{lcondi}).
\end{proof}

\end{document}
